\chapter{Heartburn}
\index{Heartburn}

\section*{Definition and Overview}
Heartburn is a burning pain or discomfort in the chest, usually behind the breastbone, that rises towards the throat. It is the most common symptom of gastro-oesophageal reflux disease (GORD), in which stomach acid flows back up into the oesophagus. Despite its name, heartburn has nothing to do with the heart. It is extremely common, usually triggered by meals, certain foods, posture, or lifestyle factors, and for most people it is a nuisance rather than a serious problem. Persistent or severe heartburn warrants assessment to exclude complications and other conditions.

\section*{Epidemiology}
Heartburn affects around one in five people at least weekly, and up to half of the population experiences it at some time. Its prevalence increases with age, and it is more common in people who are overweight, pregnant, or taking certain medications. Most people with heartburn manage it with lifestyle measures and over-the-counter medicines, while a smaller proportion have GORD requiring long-term treatment. The distinction between occasional heartburn and true GORD is based on the frequency and severity of symptoms and the presence of complications.

\section*{Aetiology and Pathophysiology}
\begin{itemize}
    \item \textbf{GORD:} reflux of stomach acid and contents into the oesophagus through a weakened lower oesophageal sphincter
    \item \textbf{Dietary triggers:} fatty foods, chocolate, caffeine, alcohol, spicy foods, citrus, and carbonated drinks
    \item \textbf{Lifestyle factors:} obesity, large meals, eating late at night, lying down after meals, and smoking
    \item \textbf{Pregnancy:} hormonal changes and pressure from the growing uterus relax the sphincter
    \item \textbf{Medications:} some drugs relax the sphincter or irritate the oesophagus
    \item \textbf{Hernia:} a hiatus hernia displaces the junction of stomach and oesophagus, promoting reflux
    \item \textbf{Pathophysiology:} refluxed acid irritates the oesophageal lining, causing the burning sensation; repeated exposure can damage the lining over time
\end{itemize}

\section*{Clinical Features}
\begin{itemize}
    \item Burning retrosternal pain rising towards the throat, worse after meals or when lying down
    \item Regurgitation: sour or bitter fluid and food returning to the mouth
    \item A sour taste, excessive belching, and a sensation of a lump in the throat
    \item Symptoms worse at night and when bending over
    \item Chronic cough, hoarse voice, and sore throat from reflux irritating the airways
    \item Dental erosion from acid exposure
\end{itemize}

\section*{Differential Diagnosis}
\begin{itemize}
    \item \textbf{Cardiac chest pain:} angina and heart attack must always be considered in chest pain, especially with exertion or risk factors
    \item Oesophageal disorders: spasm, achalasia, and oesophagitis
    \item Peptic ulcer disease and gastritis
    \item Gallbladder disease, which causes upper abdominal pain after fatty meals
    \item Anxiety and musculoskeletal chest wall pain
    \item Any new, severe, or exertional chest pain requires urgent cardiac assessment
\end{itemize}

\section*{When to Seek Medical Care}
See a doctor if heartburn is frequent (more than twice weekly), severe, wakes you at night, persists despite simple treatment, or is accompanied by difficulty swallowing, weight loss, or anaemia.

\textbf{Contact your local emergency services for:}
\begin{itemize}
    \item Chest pain or pressure that is severe, crushing, or brought on by exertion
    \item Chest pain with sweating, nausea, breathlessness, or pain radiating to the arm, neck, or jaw
    \item Fainting or collapse with chest pain
    \item Vomiting blood or passing black, tarry stools
    \item Difficulty swallowing or choking
\end{itemize}

\section*{Diagnosis and Investigations}
\begin{itemize}
    \item \textbf{Clinical diagnosis:} typical symptoms in a person without alarm features can be treated without testing
    \item \textbf{Trial of treatment:} response to acid suppression supports the diagnosis
    \item \textbf{Gastroscopy (endoscopy):} visualises the oesophagus to diagnose oesophagitis, Barrett's oesophagus, hiatus hernia, and other conditions
    \item \textbf{Ambulatory pH monitoring:} measures acid exposure when the diagnosis is uncertain
    \item \textbf{Oesophageal manometry:} assesses muscle function before some surgeries
    \item \textbf{Investigations for alarm symptoms:} endoscopy, biopsy, and imaging for difficulty swallowing, weight loss, or bleeding
\end{itemize}

\section*{Management}
\begin{itemize}
    \item \textbf{Lifestyle measures:} weight loss, smaller meals, avoiding triggers, not eating within 2--3 hours of bed, and elevating the head of the bed
    \item \textbf{Antacids and alginates:} quick symptom relief for occasional heartburn
    \item \textbf{H2 receptor antagonists:} reduce acid production for milder symptoms
    \item \textbf{Proton pump inhibitors:} the most effective acid suppression for frequent symptoms and GORD
    \item \textbf{Medication review:} stop or adjust drugs that worsen reflux where possible
    \item \textbf{Surgery:} laparoscopic fundoplication strengthens the sphincter for selected patients who cannot tolerate or do not respond to medication
    \item \textbf{Long-term review:} ongoing treatment is titrated to the lowest effective dose
\end{itemize}

\section*{Complications}
\begin{itemize}
    \item Oesophagitis: inflammation, ulceration, and bleeding of the oesophageal lining
    \item Strictures: scarring narrows the oesophagus, causing difficulty swallowing
    \item Barrett's oesophagus: replacement of the lining by intestinal-type cells, increasing oesophageal cancer risk
    \item Aspiration into the airways, causing pneumonia, chronic cough, and asthma worsening
    \item Dental erosion and gum disease
    \item Sleep disturbance and reduced quality of life
\end{itemize}

\section*{Prognosis and Outcomes}
Most people with heartburn achieve good control with lifestyle measures and acid-suppressing medicines, and occasional heartburn resolves without treatment. GORD is a chronic condition for many people, requiring maintenance therapy, but complications are uncommon when symptoms are managed well. Endoscopy allows early detection of Barrett's oesophagus and other complications, which are then monitored to reduce cancer risk. Surgery offers long-term control for selected patients, though symptoms can recur over time.

\section*{Prevention and Self-Care}
\begin{itemize}
    \item Maintain a healthy weight and avoid tight clothing around the waist
    \item Eat smaller meals and avoid eating within 2--3 hours of bedtime
    \item Identify and avoid personal trigger foods and drinks
    \item Limit alcohol, caffeine, and carbonated drinks
    \item Quit smoking
    \item Elevate the head of the bed for night-time symptoms
    \item Take medicines as directed and review them regularly with your doctor
    \item Seek medical advice before self-treating persistent heartburn long term
\end{itemize}

\section*{Sources and Further Reading}
The following reputable sources provide further information for healthcare students and patients seeking reliable, current advice about heartburn.

\begin{itemize}
    \item Healthdirect. \textit{Heartburn and reflux: causes and treatment.} https://www.healthdirect.gov.au/heartburn
    \item Gastrointestinal Society. \textit{Heartburn and gastro-oesophageal reflux.} https://badgut.org/
    \item NHS. \textit{Heartburn and acid reflux.} https://www.nhs.uk/conditions/heartburn-and-acid-reflux/
    \item Mayo Clinic. \textit{Heartburn: symptoms and causes.} https://www.mayoclinic.org/
    \item MSD Manual. \textit{Gastro-oesophageal reflux disease.} https://www.msdmanuals.com/
    \item MedlinePlus. \textit{Heartburn.} https://medlineplus.gov/heartburn.html
\end{itemize}
